\begin{table*}[t]
\caption{\label{tab:yh}Tabla de caracteres para el grupo puntual icosaédrico centrosimétrico $Y_h$ ($|G|=120$). Aquí $\tau = \frac{1+\sqrt{5}}{2} \approx 1.618$ denota la razón áurea (con $\tau^2 = \tau + 1$). Las especies con subíndice $g$ (\emph{gerade}) y $u$ (\emph{ungerade}) representan paridad par e impar bajo el centro de inversión $i$, respectivamente. La última columna lista las funciones de base canónicas lineales, pseudovectoriales y cuadráticas asociadas.}
\begin{ruledtabular}
\begin{tabular}{lccccccccccl}
$Y_h$ & $E$ & $12C_5$ & $12C_5^2$ & $20C_3$ & $15C_2$ & $i$ & $12S_{10}$ & $12S_{10}^3$ & $20S_6$ & $15\sigma$ & Funciones de base \\
\hline
$A_g$   & 1 &  1      &  1      &  1  &  1  &  1 &  1       &  1       &  1  &  1 & $x^2+y^2+z^2$ \\
$A_u$   & 1 &  1      &  1      &  1  &  1  & -1 & -1       & -1       & -1  & -1 & \\
$T_{1g}$& 3 &  $\tau$ & $1-\tau$&  0  & -1  &  3 &  $\tau$  & $1-\tau$ &  0  & -1 & $(R_x, R_y, R_z)$ \\
$T_{1u}$& 3 &  $\tau$ & $1-\tau$&  0  & -1  & -3 & $-\tau$  & $\tau-1$ &  0  &  1 & $(x, y, z)$ \\
$T_{2g}$& 3 &$1-\tau$ &  $\tau$ &  0  & -1  &  3 &$1-\tau$  &  $\tau$  &  0  & -1 & \\
$T_{2u}$& 3 &$1-\tau$ &  $\tau$ &  0  & -1  & -3 &$\tau-1$  & $-\tau$  &  0  &  1 & \\
$G_g$   & 4 & -1      & -1      &  1  &  0  &  4 & -1       & -1       &  1  &  0 & \\
$G_u$   & 4 & -1      & -1      &  1  &  0  & -4 &  1       &  1       & -1  &  0 & \\
$H_g$   & 5 &  0      &  0      & -1  &  1  &  5 &  0       &  0       & -1  &  1 & $(2z^2-x^2-y^2,\, x^2-y^2,\, xy,\, xz,\, yz)$ \\
$H_u$   & 5 &  0      &  0      & -1  &  1  & -5 &  0       &  0       &  1  & -1 & \\
\end{tabular}
\end{ruledtabular}
\end{table*}
